\chapter*{General Conclusion}
\addcontentsline{toc}{chapter}{General Conclusion}

This thesis presented the design and deployment of Agrio, a Digital Twin-based intelligent irrigation platform for precision agriculture. The work was developed as a combined Master 2 and Engineering degree thesis, bringing together scientific analysis and practical implementation. The scientific part studied Digital Twins, IoT-based smart irrigation, machine learning irrigation decisions, and agricultural water management. The engineering part implemented an end-to-end prototype that connects physical sensing, radio-frequency communication, MQTT messaging, backend processing, database persistence, ML inference, Digital Twin simulation, user interfaces, scheduling, and pump control.

The state-of-the-art review showed that Digital Twins are increasingly studied in agriculture but remain difficult to deploy due to data heterogeneity, biological complexity, cost, connectivity, calibration, and limited long-term validation. The literature also shows that many systems focus on monitoring or decision support but do not always implement a complete sensing-to-actuation loop. In this context, Agrio contributes an integrated architecture and prototype demonstrating how a human-supervised Digital Twin irrigation system can be built.

The implemented system uses a Libelium Smart Agriculture kit, XBee communication, Raspberry Pi gateway, HiveMQ/MQTT, FastAPI backend, PostgreSQL database, Random Forest-based ML service, Digital Twin state and simulation services, Next.js dashboard, Flutter mobile application, WeMos actuator, and water pump. Sensor data are parsed, normalized, stored, and converted into a synchronized virtual zone state. The ML layer predicts irrigation need and suggested timing. The simulation layer compares future moisture evolution under different scenarios. The recommendation layer supports human validation. The scheduling and device-control layer can send commands to the pump.

The main result is a working proof of concept that goes beyond classical IoT monitoring. The system provides synchronization, prediction, simulation, recommendation, validation, scheduling, and actuator feedback. It therefore demonstrates a practical implementation of Digital Twin principles in a smart irrigation context.

Nevertheless, the work remains a prototype. It was validated in a controlled experimental environment, not over multiple seasons in a commercial farm. Its ML models require further field validation and retraining. Security, authentication, sensor calibration, agronomic modeling, remote sensing integration, and multi-zone optimization require future work. These limitations are important and should be considered as the next steps toward a production-grade agricultural Digital Twin.

Future improvements include deploying the system in a real field, collecting long-term data, calibrating soil and crop parameters, integrating Open-Meteo forecasts more deeply, adding satellite or UAV imagery, improving the water balance model, comparing Random Forest with advanced models, implementing reinforcement learning for irrigation optimization, strengthening cybersecurity, adding an emergency stop mechanism, and evaluating usability with farmers and agronomists.
